\begin{figure}[htbp]
\centering
\begin{tikzpicture}

% ============================================================================
% Panel (a): Rectangular pulse (rect function)
% ============================================================================
\begin{axis}[
    name=rect,
    width=7.5cm, height=5.5cm,
    xlabel={$t$},
    ylabel={$\rect(t/\tau)$},
    xmin=-2.5, xmax=2.5,
    ymin=-0.3, ymax=1.4,
    xtick={-2, -1, -0.5, 0, 0.5, 1, 2},
    xticklabels={$-2\tau$, $-\tau$, $-\frac{\tau}{2}$, $0$, $\frac{\tau}{2}$, $\tau$, $2\tau$},
    ytick={0, 1},
    grid=major,
    grid style={gray!30},
    title={\small (a) Rectangular Pulse},
    title style={at={(0.5,1.02)}, anchor=south},
    axis lines=middle,
    axis line style={->, >=stealth, thick},
    xlabel style={at={(axis description cs:1.02,0.22)}},
    ylabel style={at={(axis description cs:0.55,1.05)}}
]

% rect(t/tau): 1 for |t| < tau/2, 0 otherwise
\addplot[UofSCGarnet, very thick, domain=-2.5:-0.501, samples=50] {0};
\addplot[UofSCGarnet, very thick, domain=-0.5:0.5, samples=50] {1};
\addplot[UofSCGarnet, very thick, domain=0.501:2.5, samples=50] {0};

% Vertical edges
\draw[UofSCGarnet, very thick] (axis cs:-0.5,0) -- (axis cs:-0.5,1);
\draw[UofSCGarnet, very thick] (axis cs:0.5,0) -- (axis cs:0.5,1);

% Discontinuity markers
\node[circle, draw=UofSCGarnet, fill=white, minimum size=4pt, inner sep=0pt, line width=0.8pt] at (axis cs:-0.5,0) {};
\node[circle, fill=UofSCGarnet, minimum size=4pt, inner sep=0pt] at (axis cs:-0.5,1) {};
\node[circle, draw=UofSCGarnet, fill=white, minimum size=4pt, inner sep=0pt, line width=0.8pt] at (axis cs:0.5,0) {};
\node[circle, fill=UofSCGarnet, minimum size=4pt, inner sep=0pt] at (axis cs:0.5,1) {};

% Width annotation
\draw[UofSC70Black, <->, >=stealth, thick] (axis cs:-0.5,1.2) -- (axis cs:0.5,1.2);
\node[UofSC70Black, font=\footnotesize, above] at (axis cs:0,1.2) {width $= \tau$};

% Area annotation
\node[UofSCGarnet, font=\footnotesize] at (axis cs:0,0.5) {Area $= \tau$};

\end{axis}

% ============================================================================
% Panel (b): Triangular pulse (tri function)
% ============================================================================
\begin{axis}[
    name=tri,
    at={(rect.outer east)},
    anchor=outer west,
    xshift=0.6cm,
    width=7.5cm, height=5.5cm,
    xlabel={$t$},
    ylabel={$\tri(t/\tau)$},
    xmin=-2.5, xmax=2.5,
    ymin=-0.3, ymax=1.4,
    xtick={-2, -1, 0, 1, 2},
    xticklabels={$-2\tau$, $-\tau$, $0$, $\tau$, $2\tau$},
    ytick={0, 1},
    grid=major,
    grid style={gray!30},
    title={\small (b) Triangular Pulse},
    title style={at={(0.5,1.02)}, anchor=south},
    axis lines=middle,
    axis line style={->, >=stealth, thick},
    xlabel style={at={(axis description cs:1.02,0.22)}},
    ylabel style={at={(axis description cs:0.55,1.05)}}
]

% tri(t/tau): 1 - |t|/tau for |t| < tau, 0 otherwise
\addplot[UofSCAtlantic, very thick, domain=-2.5:-1.001, samples=50] {0};
\addplot[UofSCAtlantic, very thick, domain=-1:0, samples=50] {1 + x};
\addplot[UofSCAtlantic, very thick, domain=0:1, samples=50] {1 - x};
\addplot[UofSCAtlantic, very thick, domain=1.001:2.5, samples=50] {0};

% Corner markers
\node[circle, fill=UofSCAtlantic, minimum size=4pt, inner sep=0pt] at (axis cs:-1,0) {};
\node[circle, fill=UofSCAtlantic, minimum size=4pt, inner sep=0pt] at (axis cs:0,1) {};
\node[circle, fill=UofSCAtlantic, minimum size=4pt, inner sep=0pt] at (axis cs:1,0) {};

% Base width annotation
\draw[UofSC70Black, <->, >=stealth, thick] (axis cs:-1,-0.15) -- (axis cs:1,-0.15);
\node[UofSC70Black, font=\footnotesize, below] at (axis cs:0,-0.15) {base $= 2\tau$};

% Peak annotation
\node[UofSCAtlantic, font=\footnotesize, above right] at (axis cs:0.1,1) {peak $= 1$};

% Area annotation
\node[UofSCAtlantic, font=\footnotesize] at (axis cs:0.6,0.35) {Area $= \tau$};

\end{axis}

% ============================================================================
% Panel (c): Gaussian pulse
% ============================================================================
\begin{axis}[
    name=gauss,
    at={(rect.outer south west)},
    anchor=outer north west,
    yshift=-0.6cm,
    width=7.5cm, height=5.5cm,
    xlabel={$t$},
    ylabel={$g(t)$},
    xmin=-3.5, xmax=3.5,
    ymin=-0.15, ymax=1.15,
    xtick={-3, -2, -1, 0, 1, 2, 3},
    xticklabels={$-3\sigma$, $-2\sigma$, $-\sigma$, $0$, $\sigma$, $2\sigma$, $3\sigma$},
    ytick={0, 0.5, 1},
    yticklabels={$0$, $\frac{1}{2}$, $1$},
    grid=major,
    grid style={gray!30},
    title={\small (c) Gaussian Pulse},
    title style={at={(0.5,1.02)}, anchor=south},
    axis lines=middle,
    axis line style={->, >=stealth, thick},
    xlabel style={at={(axis description cs:1.02,0.15)}},
    ylabel style={at={(axis description cs:0.55,1.05)}}
]

% Gaussian: exp(-t^2/(2*sigma^2)), normalized with sigma=1
\addplot[UofSCHorseshoe, very thick, domain=-3.5:3.5, samples=200] {exp(-x^2/2)};

% FWHM markers (at exp(-0.5*ln(2)^2) height)
% FWHM = 2*sqrt(2*ln(2))*sigma ≈ 2.355*sigma
% Half max at t = ±sqrt(2*ln(2)) ≈ ±1.177
\draw[UofSCRose, dashed, thin] (axis cs:-3.5,0.5) -- (axis cs:3.5,0.5);
\node[UofSCRose, font=\footnotesize, right] at (axis cs:2.8,0.55) {half max};

% FWHM width annotation
\draw[UofSCRose, <->, >=stealth, thick] (axis cs:-1.177,0.5) -- (axis cs:1.177,0.5);
\node[UofSCRose, font=\footnotesize, above] at (axis cs:0,0.55) {FWHM};

% Standard deviation markers
\draw[UofSC50Black, thin, dashed] (axis cs:-1,0) -- (axis cs:-1,0.607);
\draw[UofSC50Black, thin, dashed] (axis cs:1,0) -- (axis cs:1,0.607);
\node[circle, fill=UofSCHorseshoe, minimum size=4pt, inner sep=0pt] at (axis cs:-1,0.607) {};
\node[circle, fill=UofSCHorseshoe, minimum size=4pt, inner sep=0pt] at (axis cs:1,0.607) {};

% Formula
\node[UofSCHorseshoe, font=\footnotesize] at (axis cs:-2.2,0.85) {$\ee^{-t^2/2\sigma^2}$};

\end{axis}

% ============================================================================
% Panel (d): Sinc function
% ============================================================================
\begin{axis}[
    name=sinc,
    at={(gauss.outer east)},
    anchor=outer west,
    xshift=0.6cm,
    width=7.5cm, height=5.5cm,
    xlabel={$t$},
    ylabel={$\sinc(t/\tau)$},
    xmin=-4.5, xmax=4.5,
    ymin=-0.35, ymax=1.15,
    xtick={-4, -3, -2, -1, 0, 1, 2, 3, 4},
    xticklabels={$-4\tau$, $-3\tau$, $-2\tau$, $-\tau$, $0$, $\tau$, $2\tau$, $3\tau$, $4\tau$},
    ytick={-0.2, 0, 0.2, 0.4, 0.6, 0.8, 1},
    grid=major,
    grid style={gray!30},
    title={\small (d) Sinc Function},
    title style={at={(0.5,1.02)}, anchor=south},
    axis lines=middle,
    axis line style={->, >=stealth, thick},
    xlabel style={at={(axis description cs:1.02,0.32)}},
    ylabel style={at={(axis description cs:0.55,1.05)}}
]

% sinc(t) = sin(pi*t)/(pi*t), using t/tau scaling
% For numerical stability, define piecewise
\addplot[UofSCCongaree, very thick, domain=-4.5:-0.05, samples=200] {sin(deg(pi*x))/(pi*x)};
\addplot[UofSCCongaree, very thick, domain=0.05:4.5, samples=200] {sin(deg(pi*x))/(pi*x)};

% Value at t=0 (limit = 1)
\node[circle, fill=UofSCCongaree, minimum size=5pt, inner sep=0pt] at (axis cs:0,1) {};
\node[UofSCCongaree, font=\footnotesize, above right] at (axis cs:0.1,0.95) {$\sinc(0) = 1$};

% Mark zeros at integer multiples
\node[circle, fill=UofSCRose, minimum size=4pt, inner sep=0pt] at (axis cs:-4,0) {};
\node[circle, fill=UofSCRose, minimum size=4pt, inner sep=0pt] at (axis cs:-3,0) {};
\node[circle, fill=UofSCRose, minimum size=4pt, inner sep=0pt] at (axis cs:-2,0) {};
\node[circle, fill=UofSCRose, minimum size=4pt, inner sep=0pt] at (axis cs:-1,0) {};
\node[circle, fill=UofSCRose, minimum size=4pt, inner sep=0pt] at (axis cs:1,0) {};
\node[circle, fill=UofSCRose, minimum size=4pt, inner sep=0pt] at (axis cs:2,0) {};
\node[circle, fill=UofSCRose, minimum size=4pt, inner sep=0pt] at (axis cs:3,0) {};
\node[circle, fill=UofSCRose, minimum size=4pt, inner sep=0pt] at (axis cs:4,0) {};

% Zero label
\node[UofSCRose, font=\scriptsize] at (axis cs:2.5,-0.25) {zeros at $t = n\tau$};

% Main lobe annotation
\draw[UofSC70Black, <->, >=stealth] (axis cs:-1,0.85) -- (axis cs:1,0.85);
\node[UofSC70Black, font=\scriptsize, above] at (axis cs:0,0.85) {main lobe};

% Fourier transform note
\node[UofSCCongaree, font=\footnotesize, align=center] at (axis cs:-3,0.7) {$\mathcal{F}\{\rect\}$};

\end{axis}

\end{tikzpicture}
\caption{Common pulse-type signals and their characteristics. (a) \textbf{Rectangular pulse} $\rect(t/\tau)$: unity height over width $\tau$, fundamental building block with Fourier transform proportional to sinc. (b) \textbf{Triangular pulse} $\tri(t/\tau)$: linear rise and fall with base $2\tau$, equals the convolution $\rect \ast \rect$, smoother spectrum than rectangular. (c) \textbf{Gaussian pulse} $\ee^{-t^2/2\sigma^2}$: infinitely smooth with no sharp edges, characterized by standard deviation $\sigma$ and full-width-at-half-maximum (FWHM) $\approx 2.355\sigma$; its Fourier transform is also Gaussian. (d) \textbf{Sinc function} $\sinc(t/\tau) = \sin(\pi t/\tau)/(\pi t/\tau)$: the Fourier transform of the rectangular pulse, with zeros at integer multiples of $\tau$ and oscillating side lobes; appears in ideal low-pass filter impulse response and sampling theory.}
\label{fig:pulse_signals}
\end{figure}
